\begin{abstract}
Self-training for low-resource text classification is usually guarded by a confidence threshold, a second-model check and a held-out roll-back. We test how well these guards work for Amharic news classification with only 80 labelled articles, using paired experiments over \nSeeds{} seeds on MasakhaNEWS. With a raw MLP teacher and a 0.9 confidence threshold, an unprotected loop loses \cmpMainNaive{} macro-F1 relative to the seed-only teacher on every seed, although \precMainNaive{} of its accepted pseudo-labels are correct. The teacher is underconfident, the threshold admits almost only sports and health articles (about 1\% of business articles against 78\% of sports articles), and the training set drifts. The failure is conditional: a calibrated character-level teacher or a lower threshold removes it, and class-balanced selection turns it into a consistent gain of \cmpMainBalanced{}. A feature-diverse second-model verifier rarely disagrees and is usually the one that is wrong when it does, because the two models' errors are strongly correlated; development-set roll-back prevents the collapse but cannot certify small gains. Distilling a teacher trained on all labels into a student with about 320 times fewer parameters left accuracy unchanged (equivalence within $\pm0.01$ macro-F1), while distilling from the low-label teacher reduced accuracy relative to training the student on the seed labels (\dsDiffLowLinear{} for a linear student). The loop's small development gain did not appear on the test set (\mergedGainTeacherTest{}). We release code, per-run results and seed-set indices.
\end{abstract}

\noindent\textbf{Keywords:} Amharic; self-training; pseudo-labelling; calibration; class imbalance; knowledge distillation; low-resource NLP
